\documentclass{beamer}
\usepackage{graphicx}
\usepackage{xcolor}
\usepackage{tikz}

% 自定义颜色
\definecolor{purpleblue}{RGB}{80, 80, 180}
\definecolor{lightgray}{RGB}{230, 230, 230}
\definecolor{novelPurple}{RGB}{98, 54, 255} 
\definecolor{darkgray}{RGB}{50, 50, 50}

% 主题设置
\usetheme{Madrid}  
\usecolortheme{dolphin} 

% 设置字体（移除 fontspec，确保 pdfLaTeX 兼容）
\setbeamerfont{title}{size=\Large,series=\bfseries}
\setbeamerfont{frametitle}{size=\LARGE,series=\bfseries}
\setbeamercolor{frametitle}{fg=white, bg=novelPurple}
\setbeamercolor{title}{fg=white, bg=purpleblue}

% 封面信息
\title[SAT Unit 2.1.2 Expression ]{\textbf{SAT Unit 2.1.2\\ Expression}}
\author{\textbf{Jack Zeng}}
\institute{\textcolor{white}{\textbf{Novel Prep}}}
\date{\today} 

\begin{document}

% --- Title Slide ---
\begin{frame}
    \titlepage
    \vfill
    \begin{flushright}
        \textcolor{purpleblue}{\Huge \textbf{Novel Prep}}
    \end{flushright}
\end{frame}


\begin{frame}
    \begin{tikzpicture}[remember picture, overlay]

        \shade[top color=softPurple, bottom color=softGray] 
            (current page.north west) rectangle (current page.south east);

        \fill[white, opacity=0.3] (current page.center) circle (4cm);
        

        \node[align=center] at (current page.center) {
            {\Huge\bfseries\textcolor{purpleblue}{Novel Prep}} \\[0.5cm]
            {\Large\itshape\textcolor{lightText}{Excellence in SAT Prep}}
        };
        \foreach \x in {1,2,3,4,5} {
            \fill[purpleblue, opacity=0.2] (rand*10-5, rand*7-3.5) circle (0.1+\x*0.05);
        }
        ;
    \end{tikzpicture}
\end{frame}

\begin{frame}{Overview}
 \tableofcontents
    
\end{frame}


% Slide: Title Page Already Included
% Slide: Second Title Design Already Included

\begin{frame}
\section{Combine Like Terms and Simplify}
    \begin{tikzpicture}[remember picture, overlay]

        \shade[top color=softPurple, bottom color=softGray] 
            (current page.north west) rectangle (current page.south east);

        \fill[white, opacity=0.3] (current page.center) circle (4cm);
        

        \node[align=center] at (current page.center) {
            {\LARGE\bfseries\textcolor{purpleblue}{Combine Like Terms \& Simplify  }} \\[0.5cm]
            {\Large\itshape\textcolor{lightText}{Excellence in SAT Prep}}
        };
        \foreach \x in {1,2,3,4,5} {
            \fill[purpleblue, opacity=0.2] (rand*10-5, rand*7-3.5) circle (0.1+\x*0.05);
        }
        ;
    \end{tikzpicture}
\end{frame}

\begin{frame}{Combining Like Terms}
\small
\textbf{Definition:} Combine terms with the exact same variables and powers.

\vspace{0.3cm}
\textbf{Example:}
\[
2(2a^2 - 3a^2b^2 - 4b^2) - (a^2 + 5a^2b^2 - 10b^2)
\]
Distribute:
\[
= 4a^2 - 6a^2b^2 - 8b^2 - a^2 - 5a^2b^2 + 10b^2
\]
Combine like terms:
\[
= 3a^2 - 11a^2b^2 + 2b^2
\]

\end{frame}



\begin{frame}{Simplifying Fractional Expressions}
\small
\textbf{Rules for Simplifying:}
\begin{itemize}
    \item Look for common factors in numerator and denominator.
    \item Apply factoring where possible.
    \item Cancel shared factors.
\end{itemize}

\vspace{0.3cm}
\textbf{Example:}
\[
\frac{4xy(x+y)(x-y)}{6x(x-y)} = \frac{2y(x+y)}{3}
\]

\vspace{0.2cm}
\textbf{Note:} Do not cancel terms unless they are factors!
\end{frame}



\begin{frame}{Question}
Which expression is equivalent to
\[
(3x^3 + 8x^2 - 4x) + (7x^2 - 11x - 7)\ ?
\]

A) $3x^3 + x^2 - 15x - 7$ \\
B) $3x^3 + 15x^2 - 15x - 7$ \\
C) $10x^5 - 7x - 7$ \\
D) $15x^4 + 3x^3 - 15x^2 - 7$
\end{frame}


\begin{frame}{Answer}
Combine like terms:

\[
(3x^3 + 8x^2 - 4x) + (7x^2 - 11x - 7)
= 3x^3 + (8x^2 + 7x^2) + (-4x - 11x) - 7
= 3x^3 + 15x^2 - 15x - 7
\]

\textbf{Correct answer:} B
\end{frame}



\begin{frame}
\section{Expansion}
    \begin{tikzpicture}[remember picture, overlay]

        \shade[top color=softPurple, bottom color=softGray] 
            (current page.north west) rectangle (current page.south east);

        \fill[white, opacity=0.3] (current page.center) circle (4cm);
        

        \node[align=center] at (current page.center) {
            {\Huge\bfseries\textcolor{purpleblue}{ Expansion}} \\[0.5cm]
            {\Large\itshape\textcolor{lightText}{Excellence in SAT Prep}}
        };
        \foreach \x in {1,2,3,4,5} {
            \fill[purpleblue, opacity=0.2] (rand*10-5, rand*7-3.5) circle (0.1+\x*0.05);
        }
        ;
    \end{tikzpicture}
\end{frame}

\begin{frame}{Expansion Using FOIL}
\small
\textbf{FOIL Method} is a technique used to expand the product of two binomials.

\vspace{0.2cm}
\textbf{FOIL stands for:}
\begin{itemize}
    \item \textbf{F} – First: Multiply the first terms in each binomial
    \item \textbf{O} – Outer: Multiply the outer terms
    \item \textbf{I} – Inner: Multiply the inner terms
    \item \textbf{L} – Last: Multiply the last terms in each binomial
\end{itemize}

\end{frame}


\begin{frame}{Expansion Using FOIL}
    

\textbf{Example:} Expand $(x + 3)(x - 5)$

\[
\begin{aligned}
\text{First: } & x \cdot x = x^2 \\
\text{Outer: } & x \cdot (-5) = -5x \\
\text{Inner: } & 3 \cdot x = 3x \\
\text{Last: } & 3 \cdot (-5) = -15 \\
\end{aligned}
\]




\textbf{Combine like terms:}
\[
x^2 - 5x + 3x - 15 = x^2 - 2x - 15
\]


\textbf{Final Answer:} $x^2 - 2x - 15$
\end{frame}


\begin{frame}{\textbf{Question: Find the Value of } \( ab \)}
The equation below is true for all \( x \), where \( a \) and \( b \) are constants:
\[
(ax + 3)(5x^2 - bx + 4) = 20x^3 - 9x^2 - 2x + 12
\]

What is the value of \( ab \)?

\textbf{Options:}
\begin{enumerate}
    \item[A.] 18
    \item[B.] 20
    \item[C.] 24
    \item[D.] 40
\end{enumerate}
\end{frame}


\begin{frame}{\textbf{Solution: Matching Coefficients}}
We expand the left-hand side:
\[
(ax + 3)(5x^2 - bx + 4)= 5a x^3 - ab x^2 + 4a x + 15x^2 - 3b x + 12
\]

Now combine like terms:
\[
5a x^3 + (-ab + 15)x^2 + (4a - 3b)x + 12
\]

Match with the right-hand side:
\[
20x^3 - 9x^2 - 2x + 12
\]

Equating coefficients:
\[
5a = 20 \Rightarrow a = 4
\]
\[
-ab + 15 = -9 \Rightarrow -ab = -24 \Rightarrow ab = 24
\]

\textbf{Answer:} \(\boxed{24}\)
\end{frame}

\begin{frame}{Question}
The expression $5x^5 + 6x^4 - 8x^3$ can be rewritten as 
\[
(x^3 - hx^2)(5x^2 + 10x),
\]
where $h$ is a constant. What is the value of $h$?
\end{frame}

\begin{frame}{Answer}
We are given:
\[
(x^3 - hx^2)(5x^2 + 10x)
\]

Use distributive property:
\[
x^3(5x^2 + 10x) - hx^2(5x^2 + 10x) = 5x^5 + 10x^4 - 5h x^4 - 10h x^3
\]

Combine like terms:
\[
5x^5 + (10 - 5h)x^4 - 10h x^3
\]

Compare with the original expression:
\[
5x^5 + 6x^4 - 8x^3
\]

Match coefficients:
\[
10 - 5h = 6 \Rightarrow 5h = 4 \Rightarrow h = \frac{4}{5}
\]

\textbf{Correct answer:} $h = \dfrac{4}{5}$
\end{frame}

\begin{frame}
\section{Factoring}
    \begin{tikzpicture}[remember picture, overlay]

        \shade[top color=softPurple, bottom color=softGray] 
            (current page.north west) rectangle (current page.south east);

        \fill[white, opacity=0.3] (current page.center) circle (4cm);
        

        \node[align=center] at (current page.center) {
            {\Huge\bfseries\textcolor{purpleblue}{  Factoring}} \\[0.5cm]
            {\Large\itshape\textcolor{lightText}{Excellence in SAT Prep}}
        };
        \foreach \x in {1,2,3,4,5} {
            \fill[purpleblue, opacity=0.2] (rand*10-5, rand*7-3.5) circle (0.1+\x*0.05);
        }
        ;
    \end{tikzpicture}
\end{frame}
\begin{frame}{Factoring Basics}
\small
\textbf{Factoring} is the process of expressing a polynomial as a product of simpler expressions.

\vspace{0.3cm}
\textbf{Step 1: Factor out the Greatest Common Factor (GCF)}
\[
6x^2 + 9x = 3x(2x + 3)
\]

\vspace{0.3cm}
\textbf{Step 2: Factoring Trinomials when $a = 1$}

\textbf{Method:} Find two numbers that:
\begin{itemize}
    \item Multiply to $c$
    \item Add to $b$
\end{itemize}

\vspace{0.2cm}
\textbf{Example:}
\[
x^2 + 5x + 6 = (x + 2)(x + 3)
\]

\vspace{0.3cm}
\textbf{Tip:} Always factor out the GCF before factoring the trinomial.
\end{frame}


\begin{frame}{Factoring When $a \ne 1$ (Cross Method)}
\small
\textbf{When the leading coefficient $a \ne 1$}, we use the trial-and-error or \textbf{Chinese Cross Method }:

\vspace{0.3cm}
\textbf{Example:}
\[
6x^2 + 7x + 2
\Rightarrow (3x + 2)(2x + 1)
\]

\textbf{Reasoning:}
\begin{itemize}
    \item $3x \cdot 1 = 3x$
    \item $2x \cdot 2 = 4x$
    \item $3x + 4x = 7x$
\end{itemize}

\vspace{0.2cm}
\textbf{General Tip:}
\begin{itemize}
    \item Look for two pairs of factors of $a$ and $c$ whose cross-products add to $b$.
    \item Always check your result by expanding (FOIL method).
\end{itemize}
\end{frame}


\begin{frame}{\textbf{Question: Factoring a Quadratic Expression}}
Which of the following expressions is(are) a factor of 
\[
3x^2 + 20x - 63?
\]

\begin{itemize}
    \item[I.] \(x - 9\)
    \item[II.] \(3x - 7\)
\end{itemize}

\textbf{Options:}
\begin{enumerate}
    \item[A.] I only
    \item[B.] II only
    \item[C.] I and II
    \item[D.] Neither I nor II
\end{enumerate}
\end{frame}


\begin{frame}{\textbf{Solution: Use Factoring or Substitution}}

We are given the quadratic expression:
\[
3x^2 + 20x - 63
\]

Let's factor it:
\[
3x^2 + 20x - 63 = (x - 3)(3x + 21)
\]

But:
\[
(x - 9)(3x - 7) = 3x^2 - 7x - 27x + 63 = 3x^2 - 34x + 63 \neq 3x^2 + 20x - 63
\]

Let’s check factoring by grouping:
\[
3x^2 + 27x - 7x - 63 = 3x(x + 9) -7(x + 9) = (x + 9)(3x - 7)
\]

So the factors are:
\[
(x + 9)(3x - 7)
\]

\textbf{Correct Factor:} \(3x - 7\)
\vspace{1mm}
\textbf{Answer:} \(\boxed{\text{B. II only}}\)
\end{frame}





\begin{frame}{\textbf{Question: Finding the Smallest Possible Factor Constant}}
One of the factors of 
\[
2x^3 + 42x^2 + 208x
\]
is \(x + b\), where \(b\) is a positive constant.

\bigskip
What is the smallest possible value of \(b\)?
\end{frame}

\begin{frame}{\textbf{Solution: Factor by Grouping and Cross Multiplication}}
\scriptsize
We are given:
\[
2x^3 + 42x^2 + 208x
\]

First, factor out the GCF:
\[
2x(x^2 + 21x + 104)
\]

Now, factor the quadratic using cross multiplication (a.k.a. the AC method). We're looking for two numbers that multiply to:
\[
1 \cdot 104 = 104 \quad \text{and add up to } 21.
\]

The pair is \(13\) and \(8\), since:
\[
13 \cdot 8 = 104 \quad \text{and} \quad 13 + 8 = 21.
\]

So we factor the quadratic:
\[
x^2 + 21x + 104 = (x + 13)(x + 8)
\]

Therefore, the full factorization is:
\[
2x(x + 13)(x + 8)
\]

Since the question asks for the smallest possible positive value of \(b\) such that \(x + b\) is a factor, the answer is:
\[
\boxed{8}
\]
\end{frame}

\begin{frame}{\textbf{Question: Factoring with Decimal Coefficients}}
The expression 
\[
0.36x^2 + 0.63x + 1.17
\]
can be rewritten as 
\[
a(4x^2 + 7x + 13),
\]
where $a$ is a constant. What is the value of $a$?
\end{frame}

\begin{frame}{\textbf{Answer and Explanation}}
We are given:
\[
0.36x^2 + 0.63x + 1.17 = a(4x^2 + 7x + 13)
\]

Match the coefficients:
\[
a \cdot 4 = 0.36 \Rightarrow a = \frac{0.36}{4} = 0.09
\]

\textbf{Correct Answer:} $a = \boxed{0.09}$
\end{frame}




\begin{frame}{Factoring Rational Expressions (Part 1)}
\small
\textbf{Definition:} A rational expression is a fraction where both the numerator and the denominator are polynomials.

\vspace{0.3cm}
\textbf{Goal:} Simplify by factoring both the numerator and denominator, then cancelling common factors.

\vspace{0.4cm}
\textbf{Steps to Simplify:}
\begin{enumerate}
    \item \textbf{Factor} the numerator and denominator completely.
    \item \textbf{Identify common factors} in both.
    \item \textbf{Cancel out} the common factors.
    \item \textbf{State restrictions:} Any value that makes the denominator zero must be excluded from the domain.
\end{enumerate}

\vspace{0.4cm}
\textbf{Note:} Only factors (not terms) can be cancelled.

\[
\frac{x^2 - 9}{x^2 - x - 6} = \frac{(x - 3)(x + 3)}{(x - 3)(x + 2)} = \frac{x + 3}{x + 2} \quad \text{with } x \ne 3, -2
\]
\end{frame}




\begin{frame}{Factoring Rational Expressions (Part 2)}
\small
\textbf{Special Factoring Patterns to Know:}
\begin{itemize}
    \item Difference of squares: $a^2 - b^2 = (a - b)(a + b)$
    \item Perfect square trinomials: $a^2 \pm 2ab + b^2 = (a \pm b)^2$
    \item Difference and sum of cubes:
    \[
    a^3 - b^3 = (a - b)(a^2 + ab + b^2)
    \]
    \[
    a^3 + b^3 = (a + b)(a^2 - ab + b^2)
    \]
\end{itemize}

\vspace{0.3cm}
\textbf{Example:}
\[
\frac{4x^2 - 9}{2x + 3} = \frac{(2x - 3)(2x + 3)}{2x + 3} = 2x - 3 \quad \text{with } x \ne -\frac{3}{2}
\]

\vspace{0.3cm}
\textbf{Tip:}
\begin{itemize}
    \item Always check if the polynomial is factorable before attempting to cancel.
    \item Watch for common mistakes like cancelling terms instead of factors.
\end{itemize}
\end{frame}


\begin{frame}{\textbf{Question: Polynomial Division}}
\textbf{Given:}
\[
\frac{x^2 + 6x - 7}{x + 7} = ax + d
\]
The equation is true for all \( x \ne -7 \), where \( a \) and \( d \) are integers.

\textbf{What is the value of \( a + d \)?}

\begin{enumerate}
    \item[A.] $-6$
    \item[B.] $-1$
    \item[C.] $0$
    \item[D.] $1$
\end{enumerate}
\end{frame}

\begin{frame}{\textbf{Answer and Explanation}}
We are dividing a quadratic by a linear binomial:
\[
\frac{x^2 + 6x - 7}{x + 7}
\]

We use polynomial division or factoring. First, factor the numerator:
\[
x^2 + 6x - 7 = (x + 7)(x - 1)
\]

So:
\[
\frac{(x + 7)(x - 1)}{x + 7} = x - 1 \quad \text{for } x \ne -7
\]

Thus,
\[
ax + d = x - 1 \Rightarrow a = 1, \quad d = -1
\]

Therefore,
\[
a + d = 1 + (-1) = \boxed{0}
\]
\end{frame}


\begin{frame}{\textbf{Question: Simplifying Rational Expressions}}
Let:
\[
f(x) = x^3 - 9x, \quad g(x) = x^2 - 2x - 3
\]
Which of the following expressions is equivalent to:
\[
\frac{f(x)}{g(x)} \quad \text{for } x > 3
\]

\textbf{Options:}
\begin{enumerate}
    \item[A.] $\frac{1}{x + 1}$
    \item[B.] $\frac{x + 3}{x + 1}$
    \item[C.] $\frac{x(x - 3)}{x + 1}$
    \item[D.] $\frac{x(x + 3)}{x + 1}$
\end{enumerate}
\end{frame}

\begin{frame}{\textbf{Answer and Explanation}}
We start by factoring both the numerator and denominator:

\[
f(x) = x^3 - 9x = x(x^2 - 9) = x(x - 3)(x + 3)
\]
\[
g(x) = x^2 - 2x - 3 = (x - 3)(x + 1)
\]

Now simplify:
\[
\frac{f(x)}{g(x)} = \frac{x(x - 3)(x + 3)}{(x - 3)(x + 1)} = \frac{x(x + 3)}{x + 1}
\]

So the correct answer is:
\[
\boxed{\text{D. } \frac{x(x + 3)}{x + 1}}
\]
\end{frame}



\begin{frame}
\section{Fraction}
    \begin{tikzpicture}[remember picture, overlay]

        \shade[top color=softPurple, bottom color=softGray] 
            (current page.north west) rectangle (current page.south east);

        \fill[white, opacity=0.3] (current page.center) circle (4cm);
        

        \node[align=center] at (current page.center) {
            {\Huge\bfseries\textcolor{purpleblue}{Fraction}} \\[0.5cm]
            {\Large\itshape\textcolor{lightText}{Excellence in SAT Prep}}
        };
        \foreach \x in {1,2,3,4,5} {
            \fill[purpleblue, opacity=0.2] (rand*10-5, rand*7-3.5) circle (0.1+\x*0.05);
        }
        ;
    \end{tikzpicture}
\end{frame}

\begin{frame}{Adding Algebraic Fractions}
\small
\textbf{To add or subtract algebraic fractions:}
\begin{itemize}
    \item Step 1: Find a common denominator.
    \item Step 2: Rewrite each fraction with the common denominator.
    \item Step 3: Combine the numerators.
    \item Step 4: Simplify if possible.
\end{itemize}

\textbf{Example:} 
\[
\frac{1}{x} + \frac{1}{x+1}
\]

\textbf{Step 1:} Common denominator is $x(x+1)$

\[
\frac{1}{x} = \frac{x+1}{x(x+1)}, \quad 
\frac{1}{x+1} = \frac{x}{x(x+1)}
\]

\textbf{Step 2:} Add the numerators:
\[
\frac{x+1}{x(x+1)} + \frac{x}{x(x+1)} = \frac{(x+1) + x}{x(x+1)} = \frac{2x+1}{x(x+1)}
\]

\textbf{Tip:} Always check if you can factor or simplify the final expression.
\end{frame}



\begin{frame}{Flipping and Splitting Fractions}
\small
\textbf{1. Flipping Fractions (Reciprocals):}

\textbf{Rule:} When dividing by a fraction, multiply by its reciprocal.

\[
\frac{a}{b} \div \frac{c}{d} = \frac{a}{b} \cdot \frac{d}{c}
\]

\vspace{0.2cm}
\textbf{Example:}
\[
\frac{3}{4} \div \frac{2}{5} = \frac{3}{4} \cdot \frac{5}{2} = \frac{15}{8}
\]

\textbf{Tip:} Always flip the \textbf{second} fraction and multiply.

\vspace{0.4cm}
\textbf{2. Splitting Fractions:}

\textbf{Rule:} A fraction with addition or subtraction in the numerator can be split into two separate fractions:

\[
\frac{a + b}{c} = \frac{a}{c} + \frac{b}{c}, \quad \frac{a - b}{c} = \frac{a}{c} - \frac{b}{c}
\]

\textbf{But:} You \textbf{cannot} split a denominator:
\[
\frac{a}{b + c} \ne \frac{a}{b} + \frac{a}{c}
\]

\vspace{0.2cm}
\textbf{Example:}
\[
\frac{2x + 6}{x} = \frac{2x}{x} + \frac{6}{x} = 2 + \frac{6}{x}
\]

\textbf{Tip:} Only split the numerator when the denominator is a single term.
\end{frame}


\begin{frame}{Question:Rational Expressions}
\small
\textbf{Question:}

\[
\frac{2}{x - 2} + \frac{3}{x + 5} = \frac{rx + t}{(x - 2)(x + 5)}
\]

This equation is true for all $x > 2$, where $r$ and $t$ are positive constants.  
\textbf{What is the value of $rt$?}

\begin{itemize}
    \item[A.] $-20$
    \item[B.] $15$
    \item[C.] $20$
    \item[D.] $60$
\end{itemize}
\end{frame}

\begin{frame}{Solution: Rational Expressions}
\small
Start by finding a common denominator:

\[
\frac{2}{x - 2} + \frac{3}{x + 5} = \frac{2(x + 5) + 3(x - 2)}{(x - 2)(x + 5)}
\]

Now expand the numerator:

\[
2(x + 5) + 3(x - 2) = 2x + 10 + 3x - 6 = 5x + 4
\]

So the left-hand side becomes:

\[
\frac{5x + 4}{(x - 2)(x + 5)}
\]

Compare with the right-hand side:

\[
\frac{rx + t}{(x - 2)(x + 5)} \Rightarrow rx + t = 5x + 4
\]

So:
\[
r = 5,\quad t = 4
\Rightarrow rt = 5 \cdot 4 = \boxed{20}
\]

\textbf{Answer: (C) 20}
\end{frame}


\begin{frame}
\section{Isolate a Variable}
    \begin{tikzpicture}[remember picture, overlay]

        \shade[top color=softPurple, bottom color=softGray] 
            (current page.north west) rectangle (current page.south east);

        \fill[white, opacity=0.3] (current page.center) circle (4cm);
        

        \node[align=center] at (current page.center) {
            {\Huge\bfseries\textcolor{purpleblue}{Isolate a
Variable}} \\[0.5cm]
            {\LARGE\itshape\textcolor{lightText}{Excellence in SAT Prep}}
        };
        \foreach \x in {1,2,3,4,5} {
            \fill[purpleblue, opacity=0.2] (rand*10-5, rand*7-3.5) circle (0.1+\x*0.05);
        }
        ;
    \end{tikzpicture}
\end{frame}


\begin{frame}{Cross-Multiplication Technique}
\small
When you have an equation of the form:
\[
\frac{A}{B} = \frac{C}{D}
\]
You can use \textbf{cross-multiplication} to eliminate denominators:
\[
A \cdot D = B \cdot C
\]

\vspace{0.3cm}
\textbf{Example:}
\[
\frac{u - 3}{1} = \frac{6}{t - 2} \quad \Rightarrow \quad (u - 3)(t - 2) = 6
\]

\vspace{0.3cm}
\textbf{Why it works:} This is based on the property of proportions. Multiplying both sides by the least common denominator (LCD) clears all denominators in one step.
\end{frame}

\begin{frame}{Clearing Square Roots: Squaring Both Sides}
\small
If a variable is under a square root, square both sides to eliminate the radical.

\textbf{General Rule:}
\[
\sqrt{E} = F \quad \Rightarrow \quad E = F^2
\]

\vspace{0.3cm}
\textbf{Example:}
\[
2\sqrt{w + 19} = \frac{14x}{7y}
\quad \Rightarrow \quad
\sqrt{w + 19} = \frac{2x}{y}
\quad \Rightarrow \quad
w + 19 = \left( \frac{2x}{y} \right)^2
\]

\vspace{0.3cm}
\textbf{Tip:} Always isolate the radical before squaring both sides.
\end{frame}



\begin{frame}{Solving Literal Equations (Isolate a Variable)}
\small
Literal equations are formulas or equations involving multiple variables. The goal is to \textbf{solve for one variable in terms of the others}.

\vspace{0.2cm}
\textbf{Key Steps:}
\begin{itemize}
    \item Step 1: Identify the variable you want to isolate (make it the subject).
    \item Step 2: Use inverse operations (addition, subtraction, multiplication, division) to isolate that variable.
    \item Step 3: If needed, eliminate radicals or simplify fractions.
\end{itemize}



\vspace{0.3cm}
\textbf{Tips:}
\begin{itemize}
    \item Use cross-multiplication when variables appear in denominators.
    \item Clear square roots by squaring both sides of the equation.
    \item Watch for extraneous solutions if squaring or multiplying both sides.
\end{itemize}
\end{frame}

\begin{frame}{\textbf{Question: Solving for } $w$ \textbf{ in terms of } $x$ \textbf{ and } $y$}
Given the equation:
\[
\frac{14x}{7y} = 2\sqrt{w + 19}
\]
The equation relates the distinct positive real numbers $w$, $x$, and $y$.

Which equation correctly expresses $w$ in terms of $x$ and $y$?

\textbf{Options:}
\begin{enumerate}
    \item[A.] $w = \sqrt{\frac{x}{y}} - 19$
    \item[B.] $w = \sqrt{\frac{28x}{14y}} - 19$
    \item[C.] $w = \left(\frac{x}{y}\right)^2 - 19$
    \item[D.] $w = \left(\frac{28x}{14y}\right)^2 - 19$
\end{enumerate}
\end{frame}

\begin{frame}{\textbf{Answer and Explanation}}
Start with the equation:
\[
\frac{14x}{7y} = 2\sqrt{w + 19}
\Rightarrow 2\sqrt{w + 19} = 2\sqrt{w + 19}
\Rightarrow \frac{14x}{7y} = 2\sqrt{w + 19}
\Rightarrow \frac{2x}{y} = \sqrt{w + 19}
\]

Square both sides:
\[
\left(\frac{2x}{y}\right)^2 = w + 19
\Rightarrow \frac{4x^2}{y^2} = w + 19
\Rightarrow w = \frac{4x^2}{y^2} - 19
\]

Now compare this with the answer choices. Only option \textbf{C} correctly represents the structure:
\[
w = \left(\frac{x}{y}\right)^2 - 19
\]

\textbf{Correct Answer: C}
\end{frame}


\begin{frame}{\textbf{Question: Solving for } $t$ \textbf{ in terms of } $u$}
Given the equation:
\[
u - 3 = \frac{6}{t - 2}
\]
What is $t$ in terms of $u$?

\textbf{Options:}
\begin{enumerate}
    \item[A.] $t = \frac{1}{u}$
    \item[B.] $t = \frac{2u + 9}{u}$
    \item[C.] $t = \frac{1}{u - 3}$
    \item[D.] $t = \frac{2u}{u - 3}$
\end{enumerate}
\end{frame}

\begin{frame}{\textbf{Answer and Explanation}}
Starting from the equation:
\[
u - 3 = \frac{6}{t - 2}
\]

Multiply both sides by $(t - 2)$:
\[
(u - 3)(t - 2) = 6
\]

Expand the left-hand side:
\[
ut - 2u - 3t + 6 = 6
\Rightarrow ut - 3t = 2u
\Rightarrow t(u - 3) = 2u
\Rightarrow t = \frac{2u}{u - 3}
\]

\textbf{Correct Answer: D}
\end{frame}


\begin{frame}
\section{Practices}
    \begin{tikzpicture}[remember picture, overlay]

        \shade[top color=softPurple, bottom color=softGray] 
            (current page.north west) rectangle (current page.south east);

        \fill[white, opacity=0.3] (current page.center) circle (4cm);
        

        \node[align=center] at (current page.center) {
            {\Huge\bfseries\textcolor{purpleblue}{Practices}} \\[0.5cm]
            {\LARGE\itshape\textcolor{lightText}{Excellence in SAT Prep}}
        };
        \foreach \x in {1,2,3,4,5} {
            \fill[purpleblue, opacity=0.2] (rand*10-5, rand*7-3.5) circle (0.1+\x*0.05);
        }
        ;
    \end{tikzpicture}
\end{frame}

\begin{frame}{Question}
Which expression is equivalent to
\[
\frac{y - 3xy}{2x^2 - 2xy} + \frac{1 + 3y}{x - y}\ ?
\]

A) $\dfrac{2x - 3xy + 4y}{2x^2 - 2xy}$ \\
B) $\dfrac{2x + 3xy + y}{2x^2 - 2xy}$ \\
C) $\dfrac{2x + 3xy + y}{2x^3 - 4x^2y + 2xy^2}$ \\
D) $\dfrac{4y - 3xy + 1}{2x^3 - 4x^2y + 2xy^2}$
\end{frame}


\begin{frame}{Answer}
We simplify each term separately:

First term:
\[
\frac{y - 3xy}{2x^2 - 2xy} = \frac{y(1 - 3x)}{2x(x - y)}
\]

Second term:
\[
\frac{1 + 3y}{x - y}
\]

To combine these, we need a common denominator. The least common denominator (LCD) is $2x(x - y)$.

Rewrite the second term using LCD:
\[
\frac{1 + 3y}{x - y} = \frac{(1 + 3y) \cdot 2x}{2x(x - y)} = \frac{2x + 6xy}{2x(x - y)}
\]

Now add the two terms:
\[
\frac{y - 3xy + 2x + 6xy}{2x^2 - 2xy} = \frac{2x + 3xy + y}{2x^2 - 2xy}
\]

\textbf{Correct answer:} B
\end{frame}

\begin{frame}{Question}
If $\dfrac{4}{k + 2} = \dfrac{x}{3}$, where $k \ne -2$, which equation correctly expresses $k$ in terms of $x$?

\vspace{2em}

A) $k = \dfrac{12 - 2x}{x}$ \\
B) $k = \dfrac{12 + 2x}{x}$ \\
C) $k = \dfrac{x}{12 + 2x}$ \\
D) $k = 12x - 2$
\end{frame}


\begin{frame}{Answer}
We are given:
\[
\frac{4}{k + 2} = \frac{x}{3}
\]

Cross-multiply:
\[
4 \cdot 3 = x(k + 2) \Rightarrow 12 = xk + 2x
\]

Solve for $k$:
\[
12 - 2x = xk \Rightarrow k = \frac{12 - 2x}{x}
\]

\textbf{Correct answer:} A
\end{frame}

\begin{frame}{Question}
\[
T = 2P \sqrt{\frac{L}{G}}
\]

The given equation relates the positive real numbers $G$, $L$, $P$, and $T$. Which equation correctly expresses $G$ in terms of $L$, $P$, and $T$?

\begin{enumerate}
  \renewcommand{\labelenumi}{\Alph{enumi})}
  \item[A] $\dfrac{2PL}{T^2}$
  \item[B] $\dfrac{T^2}{4LP^2}$
  \item[C] $\dfrac{LT^2}{4P^2}$
  \item[D] $\dfrac{4LP^2}{T^2}$
\end{enumerate}
\end{frame}


\begin{frame}{Answer}
We start with the equation:
\[
T = 2P \sqrt{\frac{L}{G}}
\]

Divide both sides by $2P$:
\[
\frac{T}{2P} = \sqrt{\frac{L}{G}}
\]

Square both sides:
\[
\left( \frac{T}{2P} \right)^2 = \frac{L}{G}
\Rightarrow \frac{T^2}{4P^2} = \frac{L}{G}
\]

Solve for $G$:
\[
G = \frac{4P^2 L}{T^2}
\]

\textbf{Correct answer:} D
\end{frame}







\begin{frame}{\textbf{Question: Integer Relationship in Factored Form}}
The expression 
\[
4x^2 + bx - 45
\]
where $b$ is a constant, can be rewritten as 
\[
(hx + k)(x + j),
\]
where $h$, $k$, and $j$ are 、\textbf{integer constants}. Which of the following must be an integer?

\textbf{Options:}
\begin{enumerate}
    \item[A.] $\frac{b}{h}$
    \item[B.] $\frac{b}{k}$
    \item[C.] $\frac{45}{h}$
    \item[D.] $\frac{45}{k}$
\end{enumerate}
\end{frame}

\begin{frame}{\textbf{Answer: Integer Relationship in Factored Form}}
We are given:
\[
4x^2 + bx - 45 = (hx + k)(x + j)
\]
Expanding the right-hand side:
\[
hx^2 + (hj + k)x + kj
\]
By comparing coefficients:
\[
h = 4, \quad b = hj + k, \quad kj = -45
\]

Since $k$ and $j$ are integers and their product $kj = -45$, it follows that $k$ must divide $45$. Therefore:
\[
\frac{45}{k} \text{ must be an integer}.
\]

\textbf{Correct Answer:} \textbf{D. } $\frac{45}{k}$
\end{frame}














\begin{frame}{}
    \begin{tikzpicture}[remember picture, overlay]
        % 背景渐变色：蓝紫到白
        \shade[top color=novelPurple!40, bottom color=white]
            (current page.north west) rectangle (current page.south east);

        % 中心柔光
        \fill[white, opacity=0.15] (current page.center) circle (5cm);

        % 柔和光点点缀
        \foreach \x in {1,2,3,4,5} {
            \fill[novelPurple, opacity=0.12] (rand*10-5, rand*7-3.5) circle (0.1+\x*0.05);
        }
    \end{tikzpicture}

    \centering
    \vspace{5em}

    {\Huge \textbf{\textcolor{novelPurple}{Thank You!}}} \\[1em]
    {\large \textcolor{gray}{We appreciate your time and focus.}} \\[3em]

    {\LARGE \textbf{\textcolor{novelPurple}{\textsf{Novel Prep}}}}

    \vfill
\end{frame}





\end{document}